%%%%%%%%%%%%%%%%%%%%%%%%%%%%%%%%%%%%%%%%%%%%%%%%%%%%%%%%%%%%%%%%%%%%%%
% Lớp 11 | Quan hệ vuông góc trong không gian | Hình lăng trụ đứng. Hình chóp đều. Thể tích
% Dạng 1H8?7-5: Thể tích khối lăng trụ tứ giác
% Số câu: 25  (sinh tự động bởi tach_cau.py; từ nay thầy sửa trực tiếp file này)
%%%%%%%%%%%%%%%%%%%%%%%%%%%%%%%%%%%%%%%%%%%%%%%%%%%%%%%%%%%%%%%%%%%%%%

%%%%%%%%%%%%%%%% TRẮC NGHIỆM 4 PHƯƠNG ÁN %%%%%%%%%%%%%%%%

% nguồn: Ok_Toan11_TracNghiem_Lan2
\begin{ex}%[1H8N7-5]%[ma:77a8a258]
	Cho khối hộp chữ nhật có ba  kích thước $3$; $4$; $5$. Thể tích của khối hộp đã cho bằng
	\choice 
		{$10$} 
		 {$20$}
		{$12$} 
		{\True $60$} 
	 \loigiai{
		Thể tích khối hộp chữ nhật có ba kích thước $3$, $4$, $5$ là $V =  3 \cdot 4 \cdot 5 = 60$.
	 }
\end{ex}

% nguồn: Ok_Toan11_TracNghiem_Lan2
\begin{ex}%[1H8N7-5]%[ma:fdc5468d]
	Thể tích khối lăng trụ có  diện tích đáy $B$ và có chiều cao $h$ là
	\choice 
		 {\True $V = Bh$}
		 {$V = \dfrac{4}{3}Bh$}
		{$V  = \dfrac{1}{3}Bh$}
		{$V =  3Bh$}
	 \loigiai{
		Thể tích khối lăng trụ bằng  diện tích đáy nhân với chiều cao.\\
		Khi đó $V = B\cdot  h$.
	 }
\end{ex}

% nguồn: Ok_Toan11_TracNghiem_Lan2
\begin{ex}%[1H8N7-5]%[ma:82e1dbe7]
	Cho khối lăng trụ có đáy là hình vuông cạnh $a$ và chiều cao bằng $2a$. Thể tích  của khối lăng trụ đã cho bằng
	\choice 
		{$\dfrac{2}{3}a^3$} 
		{$\dfrac{4}{3}a^3$} 
		{\True $2a^3$} 
		{$4a^3$} 
	\loigiai{ 
		Diện tích đáy là $S_{\text{đáy}}  = a^2$.\\
		Chiều cao  $h=2a$.\\
		Thể tích khối lăng  trụ là
		$$V  = S_{\text{đáy}} \cdot h = a^2 \cdot 2a = 2a^3.$$
	} 
\end{ex}

% nguồn: Ok_Toan11_TracNghiem_Lan2
\begin{ex}%[1H8N7-5]%[ma:31c00466]
	Cho khối lăng trụ có đáy là hình vuông cạnh $a$ và chiều cao bằng $4a$. Thể  tích của khối lăng trụ đã cho bằng
	\choice 
		 {$16a^3$}
		 {\True $4a^3$}
		{$\dfrac{16}{3}a^3$} 
		{$\dfrac{4}{3}a^3$} 
	\loigiai{ 
		Diện tích đáy là  $S_{\text{đáy}} = a^2$.\\
		Chiều cao  $h=4a$.\\
		Thể  tích khối lăng trụ là
		$$V = S_{\text{đáy}} \cdot h = a^2 \cdot 4a  = 4a^3.$$
	} 
\end{ex}

% nguồn: Ok_Toan11_TracNghiem_Lan2
\begin{ex}%[1H8N7-5]%[ma:5af9459c]
	Cho khối hộp hình chữ nhật có ba  kích thước $2$; $4$; $6$. Thể tích của khối hộp đã cho bằng
	\choice 
		{$16$} 
		 {$12$}
		 {\True $48$}
		 {$8$}
	\loigiai{ 
		Thể tích của khối hộp chữ nhật  là $V = 2 \cdot 4 \cdot 6 = 48$.
	 }
\end{ex}

% nguồn: Ok_Toan11_TracNghiem_Lan2
\begin{ex}%[1H8N7-5]%[ma:72d47007]
	Thể tích khối  lập phương cạnh bằng $2$ bằng
	 \choice
		{$6$} 
		{\True $8$} 
		 {$4$}
		 {$2$}
	\loigiai{ 
		Thể tích  khối lập phương cạnh $2$ là $V = 2^3 = 8$.
	 }
\end{ex}

% nguồn: Ok_Toan11_TracNghiem_Lan2
\begin{ex}%[1H8H7-5]%[ma:6e507dda]
	Cho hình lăng trụ  đứng $ABCD.A'B'C'D'$ có đáy $ABCD$ là hình chữ nhật, $AB=a$, $AD=a\sqrt{2}$, $AB'=a\sqrt{5}$. Thể tích khối lăng trụ đã cho là
	\choice 
		 {$a^3\sqrt{2}$}
		{\True  $2a^3\sqrt{2}$}
		{$a^3\sqrt{10}$} 
		{$\dfrac{2a^3\sqrt{2}}{3}$} 
	 \loigiai{
	\immini{ 
		Ta có $S_{ABCD}=AB\cdot AD=a\cdot  a\sqrt{2}=a^2\sqrt{2}$.\\
		Xét $\triangle ABB'$ vuông tại  $B$ có \[BB'=\sqrt{AB'^2-AB^2}=\sqrt{\left(a\sqrt{5}\right)^2-a^2}=2a.\]
		Khi đó $V_{ABCD.A'B'C'D'}=S_{ABCD}\cdot BB'=a^2\sqrt{2}\cdot 2a=2a^3\sqrt{2}$. 
		}{ 
		\begin{tikzpicture}[scale=1,font=\footnotesize,line join=round,line cap=round,>=stealth]
				 \def\bc{5}
				 \def\ba{2}
				\def\h{3} 
				\def\gocB{30} 
				 \coordinate (B) at (0,0);
				 \coordinate (A) at (\gocB:\ba);
				\coordinate  (C) at (\bc,0);
				\coordinate (D) at ($(C)-(B)+(A)$); 
				\coordinate (A') at ($(A)+(90:\h)$); 
				\coordinate  (B') at ($(B)-(A)+(A')$);
				\coordinate (C')  at ($(C)-(A)+(A')$);
				\coordinate (D') at ($(D)-(A)+(A')$); 
				 \draw (B')--(B)--(C)--(D)--(D')--(A')--(B')--(C')--(D')
					(C)--(C'); 
				\draw[dashed]  (A')-- node [right]{$c$}
					(A)-- node [above]{$b$} 
					(D) 
					(A)-- node [left]{$a$} 
					(B); 
				\foreach \x/\g  in {A/-90,B/-90,C/-90,D/-90,A'/90,B'/90,C'/90,D'/90}
				 \fill (\x) circle (1pt)+(\g:3mm) node {$\x$};
				\draw pic[draw,angle  radius=0.4cm]{right angle=A'--A--D};
			\end{tikzpicture}
		} 
	 }
\end{ex}

% nguồn: Ok_Toan11_TracNghiem_Lan2
\begin{ex}%[1H8H7-5]%[ma:9a927024]
	Tính thể tích $V$ của khối lập phương  $ABCD.A'B'C'D'$, biết $AC' = a\sqrt{3}$.
	\choice 
		{\True $V = a^3$} 
		{$V  = \dfrac{3\sqrt{6}a^3}{4}$}
		 {$V = 3\sqrt{3}a^3$}
		{$V =  \dfrac{1}{3}a^3$}
	 \loigiai{
		\begin{center}
			\begin{tikzpicture}[scale=1,font=\footnotesize,line join=round,line cap=round,>=stealth]
					\coordinate[label=below left:$A$] (A)  at (0,0);
					\coordinate[label=below:$D$] (D) at (-50:3); 
					\coordinate[label=below:$B$] (B)  at (4,0);
					\coordinate[label=right:$C$] (C)  at ($(B)-(A)+(D)$);
					 \coordinate[label=above left:$A'$] (A') at ($(A)+(90:4)$);
					\coordinate[label=left:$D'$] (D') at ($(A')-(A)+(D)$); 
					\coordinate[label=right:$B'$] (B') at  ($(A')-(A)+(B)$);
					\coordinate[label=below right:$C'$] (C') at  ($(B')-(B)+(C)$);
					\draw (D)--(A)--(A')--(D')--(D)--(C)--(C')--(B')--(A') 
						(D')--(C'); 
					 \draw[dashed] (C')--(A)--(B)--(C)
						(B')--(B) 
						 (A)--(C);
					\foreach \diem in {A,D,C,B,A',D',C',B'} \fill (\diem)  circle (1pt);
					 \draw pic[draw,angle radius=0.3cm]{right angle=B--A--A'};
				\end{tikzpicture}
		\end{center}
		Gọi cạnh của khối lập phương là  $x$ ($x>0$).\\
		Độ dài đường  chéo của mặt đáy $ABCD$ là $AC = \sqrt{x^2+x^2} = x\sqrt{2}$.\\
		Xét tam giác  vuông $ACC'$ tại $C$, ta có
		\begin{eqnarray*}
			&&AC'^2 = AC^2 +  CC'^2\\
			&\Leftrightarrow& \left(a\sqrt{3}\right)^2 = \left(x\sqrt{2}\right)^2 +  x^2\\
			 &\Leftrightarrow&3a^2 = 2x^2 + x^2\\
			 &\Leftrightarrow& x^2 = a^2\\
			&\Rightarrow&  x=a.
		\end{eqnarray*}
		 Khi đó cạnh của khối lập phương là $a$.\\
		 Vậy thể tích khối lập phương là $V = a^3$.
	} 
\end{ex}

% nguồn: Ok_Toan11_TracNghiem_Lan2
\begin{ex}%[1H8V7-5]%[ma:1b6f0566]
	Cho khối hộp chữ nhật $ABCD.A'B'C'D'$ có đáy là hình vuông, $BD=2a$, góc giữa hai mặt phẳng $\left(A'BD\right)$ và $\left(ABCD\right)$ bằng $60^\circ $. Thể tích của khối hộp chữ nhật đã cho bằng 
	 \choice
		 {$\dfrac{2\sqrt{3}}{9}a^3$}
		{$6\sqrt{3}a^3$} 
		 {$\dfrac{2\sqrt{3}}{3}a^3$}
		{\True $2\sqrt{3}a^3$} 
	\loigiai{ 
	 \immini{Ta có $BD=\sqrt{2}AD\Rightarrow AD=\sqrt{2}a$.\\
		Nên $S_{ABCD}=(\sqrt{2}a)^2=2a^2$ và  $OA=\dfrac{1}{2}BD=a$.\\
		Gọi $O$ là trung điểm của  $DB$.\\
		Khi  đó $\heva{&AO\perp BD \\&A'O\perp BD}\Rightarrow \widehat{((A'BD);(ABCD))}=\widehat{(A'O;AO)}=\widehat{A'OA}\\
		\Rightarrow  \widehat{A'OA}=60^\circ$
		(Vì $\triangle A'AO$ vuông tại $A$ nên $\widehat{A'OA}$ là góc  nhọn).\\
		Xét $\triangle A'AO$ có $\tan  \widehat{A'OA}=\dfrac{AA'}{AO}$.\\
		Từ đó suy ra $AA'=AO\cdot  \tan \widehat{A'OA}=a\cdot \tan{60^\circ}=a\sqrt{3}$.\\
		Vậy $V_{ABCD.A'B'C'D'}=AA'\cdot S_{ABCD}=a\sqrt{3}\cdot 2a^2=2\sqrt{3}a^3$. 
		}{	\begin{tikzpicture}[declare function={r=3; gocB=-140;}]
		\path (0,0) coordinate (A) 
			(90:0.8*r) coordinate  (A')
			 (0:r) coordinate (D)
			(gocB:0.4*r) coordinate  (B)
			($(B)+(D)-(A)$) coordinate  (C)
			($(B)+(A')-(A)$)  coordinate (B')
			($(A')+(D)-(A)$) coordinate  (D')
			($(A')+(C)-(A)$) coordinate (C') 
			 ($(A)!0.5!(C)$) coordinate (O);
		\path pic["\scriptsize$60^\circ$",angle eccentricity=2,draw,fill=cyan,angle radius=7pt]{angle=  A'--O--A};
		\draw (A')--(B')--(C')--(D')--(A') 
			 (B')--(B)--(C)--(D)--(D')
			(C')--(C); 
		 \draw[dashed] (O)--(A')--(B)--(D)--(A')--(A)--(C)
			(B)--(A)--(D); 
		\foreach \t/\g in  {A/160,B/180,C/-90,D/0,O/-90,A'/90,B'/180,C'/90,D'/0}{
			 \fill[ball color=red] (\t) circle (1pt) node[shift={(\g:7pt)},font=\footnotesize]{$ \t $};
		} 
	\end{tikzpicture}
		 }
	 }
\end{ex}

% nguồn: Ok_Toan11_TracNghiem_Lan2
\begin{ex}%[1H8V7-5]%[ma:90f40903]
	 \immini[thm]{
	Cho khối lăng trụ đứng $ABCD.A'B'C'D'$  có đáy là hình thoi cạnh $a$, $BD = a\sqrt{3}$ và $AA' = 4a$. Thể tích của khối lăng trụ đã cho bằng
	\choice 
		{\True  $2\sqrt{3}a^3$}
		{$4\sqrt{3}a^3$} 
		{$\dfrac{2\sqrt{3}a^3}{3}$} 
		{$\dfrac{4\sqrt{3}a^3}{3}$} 
		 }{
		\begin{tikzpicture}[scale=1,font=\footnotesize,line join=round,line cap=round,>=stealth]
				 \coordinate[label=below left:$A$] (A) at (0,0);
				\coordinate[label=below:$D$]  (D) at (-50:3);
				\coordinate[label=below:$B$] (B) at (4,0); 
				 \coordinate[label=right:$C$] (C) at ($(B)-(A)+(D)$);
				\coordinate[label=above  left:$A'$] (A') at ($(A)+(90:4)$);
				\coordinate[label=left:$D'$]  (D') at ($(A')-(A)+(D)$);
				\coordinate[label=right:$B'$] (B') at ($(A')-(A)+(B)$); 
				\coordinate[label=below right:$C'$] (C') at ($(B')-(B)+(C)$); 
				 \draw (D)--(A)--(A')--(D')--(D)--(C)--(C')--(B')--(A')
					(D')--(C'); 
				\draw[dashed]  (A)--(B)--(C)
					 (B')--(B);
				\foreach \diem in {A,D,C,B,A',D',C',B'} \fill  (\diem) circle (1pt);
				 \draw pic[draw,angle radius=0.3cm]{right angle=B--A--A'};
			\end{tikzpicture}
		} 
	\loigiai{ 
		\begin{center}
			\begin{tikzpicture}[scale=1,font=\footnotesize,line join=round,line cap=round,>=stealth]
					\coordinate[label=below  left:$A$] (A) at (0,0);
					\coordinate[label=below:$D$]  (D) at (-50:3);
					\coordinate[label=right:$B$]  (B) at (4,0);
					\coordinate[label=right:$C$]  (C) at ($(B)-(A)+(D)$);
					 \coordinate[label=above left:$A'$] (A') at ($(A)+(90:4)$);
					 \coordinate[label=left:$D'$] (D') at ($(A')-(A)+(D)$);
					\coordinate[label=right:$B'$] (B') at ($(A')-(A)+(B)$); 
					\coordinate[label=below right:$C'$] (C')  at ($(B')-(B)+(C)$);
					\coordinate[label=below:$I$] (I) at (intersection of  A--C and B--D);
					 \draw (D)--(A)--(A')--(D')--(D)--(C)--(C')--(B')--(A')
						 (D')--(C');
					 \draw[dashed] (A)--(B)--(C)
						 (B')--(B)
						 (A)--(C)
						 (B)--(D);
					\foreach \diem in {A,D,C,B,A',D',C',B',I} \fill  (\diem) circle (1pt);
					\draw pic[draw,angle  radius=0.3cm]{right angle=B--A--A'};
				\end{tikzpicture}
		\end{center}
		Gọi $I$ là giao  điểm của hai đường chéo $AC$ và $BD$.\\
		Vì $ABCD$ là hình thoi nên $AC \perp  BD$ tại $I$ và $I$ là trung điểm của mỗi đường chéo.\\
		Ta có $IB = \dfrac{BD}{2}  = \dfrac{a\sqrt{3}}{2}$.\\
		Xét tam giác $AIB$ vuông  tại $I$, ta có
		$$AI = \sqrt{AB^2 - IB^2} = \sqrt{a^2 -  \left(\dfrac{a\sqrt{3}}{2}\right)^2} = \sqrt{a^2 - \dfrac{3a^2}{4}} = \sqrt{\dfrac{a^2}{4}} = \dfrac{a}{2}.$$
		Suy ra  $AC = 2AI = a$.\\
		Diện tích đáy hình thoi là $$S_{ABCD} = \dfrac{1}{2}\cdot AC \cdot BD = \dfrac{1}{2}  \cdot a \cdot a\sqrt{3} = \dfrac{a^2\sqrt{3}}{2}.$$
		Chiều cao lăng trụ là $h  = AA' = 4a$.\\
		Thể tích  khối lăng trụ là
		$$V = S_{ABCD} \cdot AA' = \dfrac{a^2\sqrt{3}}{2} \cdot  4a = 2\sqrt{3}a^3.$$
	} 
\end{ex}

% nguồn: Ok_Toan11_TracNghiem_Lan2
\begin{ex}%[1H8V7-5]%[ma:c9bebc15]
	Cho  khối hộp chữ nhật $ABCD.A'B'C'D'$ có đáy là hình vuông, $BD=4a$, góc giữa hai mặt phẳng $\left(A'BD\right)$ và $\left(ABCD\right)$ bằng $60^\circ $. Thể tích của khối hộp chữ nhật đã cho bằng
	 \choice
		{$48\sqrt{3}a^3$} 
		 {$\dfrac{16\sqrt{3}}{9}a^3$}
		 {$\dfrac{16\sqrt{3}}{3}a^3$}
		 {\True $16\sqrt{3}a^3$}
	 \loigiai{
	\immini{Đặt $x=AA'$,  $AB=AD=a\sqrt{8}$.\\
		Ta  có $\dfrac{\sqrt{3}}{2}=\sin \left(\left(A'BD\right),\left(ABD\right)\right)=\dfrac{d\left(A,\left(A'BD\right)\right)}{d\left(A,BD\right)}$.\\
		 Suy ra $d\left(A,\left(A'BD\right)\right)=\dfrac{\sqrt{3}}{2}\cdot 2a=a\sqrt{3}$.\\
		Vì $ABDA'$ là tam diện vuông tại $A$ nên ta  có
		\begin{center}
			$\dfrac{1}{3a^2}=\dfrac{1}{8a^2}+\dfrac{1}{8a^2}+\dfrac{1}{x^2}\Leftrightarrow x=a\sqrt{12}$. 
		\end{center}
		 Vậy $V_{ABCD.A'B'C'D'}=a\sqrt{12}\cdot a\sqrt{8}\cdot a\sqrt{8}=16\sqrt{3}a^3$.
		}{	\begin{tikzpicture}[declare function={r=3; gocB=-140;}]
		 \path (0,0) coordinate (A)
			(90:0.8*r) coordinate (A') 
			 (0:r) coordinate (D)
			 (gocB:0.4*r) coordinate (B)
			($(B)+(D)-(A)$) coordinate (C) 
			 ($(B)+(A')-(A)$) coordinate (B')
			($(A')+(D)-(A)$) coordinate (D') 
			($(A')+(C)-(A)$)  coordinate (C')
			($(A)!0.5!(C)$)  coordinate (O);
		\path pic["\scriptsize$60^\circ$",angle eccentricity=2,draw,fill=cyan,angle  radius=7pt]{angle= A'--O--A};
		\draw  (A')--(B')--(C')--(D')--(A')
			(B')--(B)--(C)--(D)--(D') 
			(C')--(C); 
		\draw[dashed] (O)--(A')--(B)--(D)--(A')--(A)--(C) 
			(B)--(A)--(D); 
		\foreach \t/\g  in {A/160,B/180,C/-90,D/0,O/-90,A'/90,B'/180,C'/90,D'/0}{
			\fill[ball color=red] (\t) circle  (1pt) node[shift={(\g:7pt)},font=\footnotesize]{$ \t $};
		 }
	\end{tikzpicture}
		} 
	} 
\end{ex}

% nguồn: Ok_Toan11_TracNghiem_Lan2
\begin{ex}%[1H8V7-5]%[ma:f35593c4]
	Cho khối hộp chữ nhật $ABCD.A'B'C'D'$có đáy là hình vuông, $BD=4a$, góc giữa 2  mặt phẳng $\left(A'BD\right),\left(ABCD\right)$ bằng $30^\circ $. Thể tích của khối hộp chữ nhật đã cho bằng:
	 \choice
		{$\dfrac{16\sqrt{3}}{9}a^3$} 
		 {$48\sqrt{3}a^3$}
		 {\True $\dfrac{16\sqrt{3}}{3}a^3$}
		 {$16\sqrt{3}a^3$}
	\loigiai{ 
	\immini{Gọi $O$ là tâm của hình vuông $ABCD$, từ giả thiết ta  có\\ $AC=4a$, $AB=\dfrac{4a}{\sqrt{2}}=2a\sqrt{2}\Rightarrow AO=2a$.\\
		Suy ra $S_{ABCD}=\left(2a\sqrt{2}\right)^2=8a^2$.\\ 
		 Mà $ABCD$ là hình vuông $\Rightarrow AO\perp BD$.\\
		 Ta có
		$\heva{& AO\perp BD \\& AA'\perp BD} \Rightarrow BD\perp  \left(A'AO\right)\Rightarrow BD\perp A'O.\\
		 \Rightarrow \left(\left(A'BD\right),\left(ABCD\right)\right)=\widehat{A'OA}$
		(tam giác $A'OA$  vuông tại $A$).
		}{	\begin{tikzpicture}[declare function={r=3; gocB=-140;}]
		\path (0,0)  coordinate (A)
			(90:0.8*r) coordinate  (A')
			(0:r)  coordinate (D)
			(gocB:0.4*r)  coordinate (B)
			($(B)+(D)-(A)$)  coordinate (C)
			($(B)+(A')-(A)$) coordinate (B') 
			($(A')+(D)-(A)$) coordinate (D') 
			($(A')+(C)-(A)$)  coordinate (C')
			($(A)!0.5!(C)$) coordinate  (O);
		\path pic["\scriptsize$30^\circ$",angle eccentricity=2,draw,fill=cyan,angle radius=7pt]{angle= A'--O--A}; 
		\draw  (A')--(B')--(C')--(D')--(A')
			(B')--(B)--(C)--(D)--(D') 
			 (C')--(C);
		\draw[dashed] (O)--(A')--(B)--(D)--(A')--(A)--(C) 
			 (B)--(A)--(D);
		\foreach \t/\g  in {A/160,B/180,C/-90,D/0,O/-90,A'/90,B'/180,C'/90,D'/0}{
			\fill[ball color=red] (\t) circle  (1pt) node[shift={(\g:7pt)},font=\footnotesize]{$ \t $};
		} 
	\end{tikzpicture}
		}\noindent 
		Từ giả thiết $\Rightarrow \widehat{A'OA}=30^\circ \Rightarrow  \tan 30^\circ =\dfrac{A'A}{AO}\Rightarrow A'A=\dfrac{\sqrt{3}}{3}\cdot 2a=\dfrac{2a\sqrt{3}}{3}$.\\
		$\Rightarrow {V_{ABCD.A'B'C'D'}}=A'A\cdot S_{ABCD}=\dfrac{2a\sqrt{3}}{3}\cdot 8a^2=\dfrac{16\sqrt{3}a^3}{3}$. 
	 }
\end{ex}

% nguồn: Ok_Toan11_TracNghiem_Lan2
\begin{ex}%[1H8V7-5]%[ma:316a3d68]
	Cho khối hộp chữ  nhật $ABCD.A'B'C'D'$ có đáy là hình vuông, $BD=2a$, góc giữa hai mặt phẳng $\left(A'BD\right)$ và $\left(ABCD\right)$ bằng $30^\circ$. Thể tích của khối hộp chữ nhật đã cho bằng
	\choice 
		{$6\sqrt{3}a^3$} 
		 {$\dfrac{2\sqrt{3}}{9}a^3$}
		{$2\sqrt{3}a^3$} 
		 {\True $\dfrac{2\sqrt{3}}{3}a^3$}
	\loigiai{ 
	\immini{Gọi $\varphi $  là góc giữa hai mặt phẳng $\left(A'BD\right)$ và $\left(ABCD\right)$.\\
		 Gọi $O=AC\cap BD$.\\
		Có $\heva{&AO\perp  BD\\&AA'\perp BD}\Rightarrow A'O\perp BD\Rightarrow \varphi =(AO;A'O)=\widehat{AOA'}=30^\circ$.\\
		Vì đáy $ABCD$ là hình vuông có $BD=2a\Rightarrow AB=AD=a\sqrt{2}$.\\ 
		Ta  có $AO=\dfrac{1}{2}AC=\dfrac{1}{2}BD=a$.\\
		Trong  $\triangle AOA'$ có $AA'=AO\cdot \tan{30^\circ}=\dfrac{a\sqrt{3}}{3}$.
		}{	\begin{tikzpicture}[declare function={r=3; gocB=-140;}]
		\path (0,0)  coordinate (A)
			(90:0.8*r) coordinate (A') 
			(0:r) coordinate  (D)
			(gocB:0.4*r) coordinate  (B)
			($(B)+(D)-(A)$)  coordinate (C)
			($(B)+(A')-(A)$)  coordinate (B')
			 ($(A')+(D)-(A)$) coordinate (D')
			($(A')+(C)-(A)$)  coordinate (C')
			($(A)!0.5!(C)$)  coordinate (O);
		\path pic["\scriptsize$30^\circ$",angle eccentricity=2,draw,fill=cyan,angle radius=7pt]{angle=  A'--O--A};
		 \draw (A')--(B')--(C')--(D')--(A')
			 (B')--(B)--(C)--(D)--(D')
			(C')--(C); 
		\draw[dashed] (O)--(A')--(B)--(D)--(A')--(A)--(C) 
			 (B)--(A)--(D);
		\foreach \t/\g in  {A/160,B/180,C/-90,D/0,O/-90,A'/90,B'/180,C'/90,D'/0}{
			\fill[ball color=red] (\t) circle (1pt) node[shift={(\g:7pt)},font=\footnotesize]{$ \t $}; 
		} 
	\end{tikzpicture}
		 }\noindent
		Vậy thể tích khối hộp chữ nhật $ABCD.A'B'C'D'$  là
		\begin{center}
			$V_{ABCD.A'B'C'D'}=AA'\cdot  S_{ABCD}=\dfrac{a\sqrt{3}}{3}\cdot 2a^2=\dfrac{2\sqrt{3}a^3}{3}$.
		\end{center}
	} 
\end{ex}


%%%%%%%%%%%%%%%% TRẢ LỜI NGẮN %%%%%%%%%%%%%%%%

% nguồn: Ok_TraLoiNgan_Lop11_Lan2
\begin{ex}%[1H8H7-5]%[ma:30c3a85f]
	\immini{Một bể cá được làm bằng kính có dạng hình hộp chữ nhật có ba kích thước là $0{,}6$ m; $2$ m; $0{,}8$ m. Tính thể tích của bể cá đó (\textit{đơn vị mét khối, làm tròn kết quả đến hàng phần trăm}).}{\begin{tikzpicture}[declare function={a=4; b=0.45*a; h=0.65*a;}]
		\path  (0:0) coordinate (A)
			(220:b) coordinate  (B)
			($(B)+(a,0)$) coordinate  (C)
			(a,0) coordinate  (D)
		 \foreach \x in {A,B,C,D}{(\x)++(90:h) coordinate (\x1)}
			($(B)!0.5!(C)$) coordinate  (E) node[below]{$2$ m}
			($(C)!0.5!(D)$) coordinate (F) node[below right]{$0{,}6$  m}
			($(B)!0.5!(B1)$) coordinate (F) node[left]{$0{,}8$  m};
		\draw[dashed] (A1)--(A)--(B) 
			(A)--(D); 
		\draw (B1)--(A1)--(D1)--(C1) 
			(D1)--(D)--(C) 
			 (B)--(C)--(C1)--(B1)--cycle;
	\end{tikzpicture}
	 }
	\shortans{0,96} 
	 \loigiai{
		Thể tích của bể cá hình hộp chữ nhật là 
		$V=abc=0{,}6\cdot 2\cdot 0{,}8=0{,}96$  (m$^3$).
	} 
\end{ex}

% nguồn: Ok_TraLoiNgan_Lop11_Lan1
\begin{ex}%[1H8V7-5]%[ma:32a78fc8]
	\immini 
	{ 
	Một cái hộp hình lập phương, bên trong nó đựng một mô hình đồ chơi có dạng hình chóp tứ giác đều mà đỉnh của hình chóp đó trùng với tâm của một mặt chiếc hộp, giả sử hình vuông đáy của hình chóp trùng với một mặt của chiếc hộp (mặt này cùng với mặt chứa đỉnh hình chóp là hai mặt đối nhau). Biết cạnh của chiếc hộp  bằng ${30 {~cm}}$, hãy tính thể tích phần không gian bên trong chiếc hộp không bị chiếm bởi mô hình đồ chơi dạng hình chóp (mô hình đồ chơi được làm bởi chất liệu nhựa đặc bên trong, kết quả có đơn vị là dm$^3$).
	}{ 
	\begin{tikzpicture}[scale=1,font=\footnotesize,line join=round,line cap=round,>=stealth]
			\path 
				(0,0) coordinate (A) 
				 ++(-130:2) coordinate (B)
				++(0:4)  coordinate (C)
				($(A)+(C)-(B)$) coordinate (D); 
			\foreach \i in {A,B,C,D}{ 
				\coordinate (\i') at  ($(\i)+(0,4)$);
			 }
			 \coordinate (O) at ($(A')!1/2!(C')$);
			 \draw (A')--(B')--(C')--(D')--cycle;
			 \draw (B)--(B')
				 (C)--(C')
				(D)--(D') 
				(B)--(C)--(D); 
			 \draw[dashed,thin](B)--(A)--(A')
				(A)--(D) 
				 (O)--(A)
				(O)--(B) 
				(O)--(C) 
				 (O)--(D);
			\foreach  \i/\g in {A'/90,B'/90,C'/90,D'/90,A/-90,B/-90,C/-90,D/-90}
			\fill[black] (\i) circle  (1pt)+(\g:5mm);
		\end{tikzpicture}
	 }
	\shortans{18} 
	 \loigiai{
		Thể tích cái  hộp (khối lập phương) là: $V_1=30^3=27000$ (cm$^3$).\\
		Xét đồ chơi có dạng hình chóp tứ giác đều, chiều cao của hình chóp bằng với một cạnh của hình lập phương, hay $h=30$ cm, đáy  của hình chóp có diện tích $S=30^2=900$ (cm$^2$).\\
		Thể tích khối  đồ chơi (khối chóp tứ giác đều) là:
		$V_2=\dfrac{1}{3} S h=\dfrac{1}{3} \cdot 900  \cdot 30=9000$ (cm$^3$).\\
		Thể tích phần không gian bên trong chiếc hộp không bị chiếm bởi mô hình đồ chơi dạng hình chóp: $V=V_1-V_2=27000-9000=18000$ (cm$^3$)=18  (dm$^3$).
	 }
\end{ex}

% nguồn: Ok_TraLoiNgan_Lop11_Lan2
\begin{ex}%[1H8V7-5]%[ma:8f3b204a]
	\immini{Một hình hộp chữ nhật $ABCD.A'B'C'D'$ có ba kích thước là $2$ cm, $3$ cm và $6$ cm. Tính thể tích của khối tứ diện $ACB'D'$.}{\begin{tikzpicture}[declare function={a=3; b=0.45*a; h=0.65*a;}]
		 \path (0:0) coordinate (A)
			(220:b) coordinate (B) 
			($(B)+(a,0)$)  coordinate (C)
			 (a,0) coordinate (D)
		\foreach \x in  {A,B,C,D}{(\x)++(90:h) coordinate (\x1)};
		 \draw[dashed] (A1)--(A)--(B)
			 (A)--(D)
			 (A)--(B1)
			(A)--(D1) 
			 (A)--(C);
		\draw (B1)--(C) 
			 (D1)--(C)
			(B1)--(D1) 
			(B1)--(A1)--(D1)--(C1) 
			(D1)--(D)--(C) 
			(B)--(C)--(C1)--(B1)--cycle; 
		\foreach \x/\goc/\t in  {A/-90/A,B/-90/D,C/-90/C,D/0/B,A1/90/A',B1/180/D',C1/0/C',D1/0/B'}
		\draw[fill=white] (\x)  circle (1pt) node[shift={(\goc:7pt)},font=\scriptsize]{$\t$};
	\end{tikzpicture}
	 }
	\shortans{12} 
	\loigiai{ 
		\begin{eqnarray*}
			V_{ACB'D'}&=&V_{ABCD.  A'B'C'D'}-(V_{D'.ACD}+V_{B'.ABC}+V_{A.ABD}+V_{C.BCD})\\
			&=&2\cdot3\cdot6-4\cdot  V_{D'.ACD}\\
			&=&36-4\cdot\dfrac{1}{3}\cdot\dfrac{1}{2}\cdot 2\cdot 3\cdot  6\\
			&=&12\,\left(\mathrm{cm}^3\right). 
		\end{eqnarray*}
	 }
\end{ex}


%%%%%%%%%%%%%%%% TỰ LUẬN %%%%%%%%%%%%%%%%

% nguồn: Ok_TuanLuan_Toan11_Lan2
\begin{ex}%[1H8H7-5]%[ma:c9ccabb7]
	Cho khối chóp $S.ABCD$ có đáy là hình vuông  cạnh $a$, mặt bên $SAB$ là tam giác cân tại $S$ và nằm trong mặt phẳng vuông góc với mặt phẳng đáy, $SC$ tạo với $(SAB)$ một góc $45^\circ$. Tính thể tích của khối chóp đã cho.
	\loigiai{ 
		\begin{center}
			\begin{tikzpicture}[scale=0.8,font=\footnotesize,line join=round,line cap=round,>=stealth]
					 \path
						(0,0) coordinate (A) 
						 (5,0) coordinate (D)
						(2,-1.5) coordinate  (C)
						($(A)+(C)-(D)$)  coordinate (B)
						($(A)!0.5!(B)$) coordinate  (H)
						($(H)+(0,5)$) coordinate (S); 
					 \draw (S)--(B)--(C)--(D)--(S)
						(S)--(C); 
					\draw[dashed](S)--(A)--(B) 
						 (A)--(D)
						(S)--(H); 
					\foreach \d/\g in {A/180,B/-90,C/-90,S/90,H/180,D/0}{ 
						\draw[fill](\d) circle (1pt) +(\g:.35)  node{$\d$};}
					\draw  pic[draw,blue,angle radius=4mm] {angle = B--S--C};
				\end{tikzpicture}
		\end{center}
		Do $\triangle SAB$ cân tại $S$ và $(SAB)\perp(ABCD)$ nên hình chiếu  vuông góc của $S$ trên $(ABCD)$ là trung điểm $H$ của $AB\Rightarrow SH \perp(ABCD)$.\\
		Ta có $BC\perp AB$, $BC\perp  SH\Rightarrow BC \perp(SAB) \Rightarrow (SC,(SAB))=\widehat{BSC}=45^\circ\Rightarrow SB=BC=a$.\\
		Suy ra $\triangle SAB$ đều cạnh $a \Rightarrow SH=\dfrac{\sqrt{3}a}{2}$.\\ 
		 Ta có $
		S_{ABCD}=a^2 \Rightarrow V_{ABCD}=\dfrac{1}{3}SH\cdot S_{ABCD}=\dfrac{\sqrt{3}a^3}{6}$. 
	 }
\end{ex}

% nguồn: Ok_TuanLuan_Toan11_Lan2
\begin{ex}%[1H8H7-5]%[ma:398e544d]
	Cho khối hộp chữ nhật $ABCD.A'B'C'D'$ có $AB=3$, $AC=5$, $AA'=8$. Thể tích  khối hộp đã cho bằng
	\loigiai{ 
	\immini 
		 {
		 Ta có $AD=BC=\sqrt{AC^2-AB^2}=4$.\\
		Thể tích khối hộp  chữ nhật đã cho là $V=AB\cdot AD\cdot AA'=3\cdot 4 \cdot 8=96$.
		} 
		 {
		\begin{tikzpicture}[font=\footnotesize,line cap=round,line join=round,scale=.7,>=stealth]
				 \def\ab{3}
				\def\h{2.7} 
				 \path
					(0,0)  coordinate (A)
					(\ab,0) coordinate (B) 
					 (-60:2) coordinate (D)
					 ($(B)+(D)$) coordinate (C);
				 \foreach \i in {A,B,C,D} \path
					 ($(\i)+(90:\h)$) coordinate (\i');
				\draw (A)--(D)--(C) 
					 (A')--(D')--(C')--(B')--cycle (A)--(A')
					(C)--(C') 
					 (D)--(D');
				 \draw[dashed](A)--(B)
					(C)--(A) 
					(C)--(B)--(B'); 
				 \foreach \p/\r in {A/180,B/0,C/-45,D/-135,A'/180,B'/0,C'/-45,D'/-135}
				\fill (\p) circle  (1pt) node[shift={(\r:3mm)}]{$\p$};
			\end{tikzpicture}
		} 
	} 
\end{ex}

% nguồn: Ok_TuanLuan_Toan11_Lan2
\begin{ex}%[1H8H7-5]%[ma:3c0e4471]
	Khối lập phương $ABCD.A'B'C'D'$ có độ dài đoạn $A'C=a$. Thể  tích của khối đó là
	\loigiai{ 
	 \immini
		 {
		Gọi độ dài  cạnh hình lập phương là $x$ với $x>0$.\\
		Ta có  $A'C=\sqrt{A'A^2+AC^2}=x\sqrt{3}=a\Leftrightarrow x=\dfrac{a\sqrt{3}}{3}$.\\
		 Vậy thể tích khối lập phương $V=x^3=\dfrac{a^3\sqrt{3}}{9}$.
		} 
		{ 
		\begin{tikzpicture}[font=\footnotesize,line cap=round,line join=round,scale=.9,>=stealth]
				 \def\ab{3}
				\def\h{2.7} 
				 \path
					 (0,0) coordinate (A)
					(\ab,0)  coordinate (B)
					(-60:2) coordinate  (D)
					($(B)+(D)$) coordinate  (C);
				\foreach  \i in {A,B,C,D} \path
					($(\i)+(90:\h)$) coordinate  (\i');
				 \draw (A)--(D)--(C)
					(A')--(D')--(C')--(B')--cycle  (A)--(A')
					 (C)--(C')
					(D)--(D'); 
				 \draw[dashed](A)--(B)
					 (A')--(C)--(A)
					(C)--(B)--(B'); 
				\foreach \p/\r in {A/180,B/0,C/-45,D/-135,A'/180,B'/0,C'/-45,D'/-135} 
				\fill (\p) circle (1pt) node[shift={(\r:3mm)}]{$\p$}; 
			\end{tikzpicture}
		 }
	 }
\end{ex}

% nguồn: Ok_TuanLuan_Toan11_Lan2
\begin{ex}%[1H8H7-5]%[ma:b3803ad7]
	Cho hình hộp đứng $ABCD.A'B'C'D'$ có đáy $ABCD$ là hình thoi cạnh $a$, góc $\widehat{BAD}=60^{\circ}$. Biết rằng góc giữa đường chéo $BD'$ và mặt đáy là $60^{\circ}$.  Thể tích khối hộp đã cho là
	 \loigiai{
	\immini{Do $ABCD.A'B'C'D'$ là hình hộp đứng nên $DD'\perp (ABCD)$,  suy ra. Do đó $(BD',(ABCD))=\widehat{DBD'}=60^{\circ}$.\\
		Do $ABCD$ là hình thoi cạnh $a$ và $\widehat{BAD}=60^{\circ}$ nên $BD=a$ và $AC=a\sqrt{3}$.\\ Suy ra $DD'=a\sqrt{3}$  và $S_{ABCD}=\dfrac{a^2\sqrt{3}}{2}$
		Vậy thể  tích khối hộp là $V=DD'\cdot S_{ABCD}=\dfrac{3}{2}a^3$.
		} 
		{\begin{tikzpicture}[scale=0.8,font=\footnotesize,line join=round,line cap=round,>=stealth]
		\def\bc{4} 
		 \def\ba{2}
		 \def\h{4}
		 \def\gocB{35}
		\coordinate[label=below left:$B$]  (B) at (0,0);
		\coordinate[label=above left:$A$]  (A) at (\gocB:\ba);
		\coordinate[label=below:$C$] (C)  at (\bc,0);
		\coordinate[label=right:$D$] (D) at  ($(C)-(B)+(A)$);
		\coordinate[label=above left:$A'$] (A')  at ($(A)+(90:\h)$);
		\coordinate[label=left:$B'$] (B')  at ($(B)-(A)+(A')$);
		 \coordinate[label=below right:$C'$] (C') at ($(C)-(A)+(A')$);
		\coordinate[label=right:$D'$] (D')  at ($(D)-(A)+(A')$);
		\draw (B')--(B)--(C)--(D)--(D')--(A')--(B')--(C')--(D') 
			 (C)--(C');
		\draw[dashed] (A')--(A)--(D)--(B) 
			 (A)--(B)--(D');
		\foreach \diem in {A,B,C,D,A',B',C',D'} \fill (\diem) circle (1pt); 
	\end{tikzpicture}
		} 
	} 
\end{ex}

% nguồn: Ok_TuanLuan_Toan11_Lan2
\begin{ex}%[1H8H7-5]%[ma:e6d31035]
	Cho hình hộp đứng $ABCD.A'B'C'D'$ có đáy là hình vuông  cạnh $a$. Khoảng cách từ điểm $A$ đến mặt phẳng $(A'BCD')$ bằng $\dfrac{a\sqrt{3}}{2}$. Tính thể tích khối hộp theo $a$.
	 \loigiai{
		\begin{center}
			\begin{tikzpicture}[scale=0.7,font=\footnotesize,line join=round,line cap=round,>=stealth]
					 \def\bc{4}
					\def\ba{2} 
					\def\h{4} 
					\def\gocB{35} 
					\coordinate[label=below left:$B$]  (B) at (0,0);
					\coordinate[label=below:$A$]  (A) at (\gocB:\ba);
					\coordinate[label=below:$C$] (C) at  (\bc,0);
					 \coordinate[label=right:$D$] (D) at ($(C)-(B)+(A)$);
					\coordinate[label=above left:$A'$]  (A') at ($(A)+(90:\h)$);
					\coordinate[label=left:$B'$] (B')  at ($(B)-(A)+(A')$);
					\coordinate[label=below right:$C'$] (C') at ($(C)-(A)+(A')$); 
					 \coordinate[label=right:$D'$] (D') at ($(D)-(A)+(A')$);
					\coordinate[label=right:$H$]  (H) at ($(A')!1/2!(B)$);
					 \draw (B')--(B)--(C)--(D)--(D')--(A')--(B')--(C')--(D')
						(C)--(C') 
						(D')--(C); 
					 \draw[dashed] (A')--(A)--(D)
						(A)--(B) 
						 (A')--(B)
						 (A)--(H);
					\foreach \diem  in {A,B,C,D,A',B',C',D',H} \fill (\diem) circle (1pt);
				\end{tikzpicture}
		\end{center}
		 Gọi $H$ là hình chiếu của $A$ lên cạnh $A'B$ \\
		$ \Rightarrow  AH\perp A'BCD'\Rightarrow AH=\dfrac{a\sqrt{3}}{2} $.\\
		Gọi $AA'=x>0$. Áp dụng hệ thức về cạnh và đường cao trong tam  giác $AA'B$.\\
		$$\dfrac{1}{AH^2}=\dfrac{1}{AA'^2}+\dfrac{1}{AB^2}\Leftrightarrow\dfrac{4}{3a^2}=\dfrac{1}{x^2}+\dfrac{1}{a^2}\Leftrightarrow  x^2=3a^2\Leftrightarrow x=a\sqrt{3}.$$
		 Khi đó
	$V_{ABCD.A'B'C'D'}=AA'\cdot AB\cdot AD=a\sqrt{3}\cdot a\cdot  a=a^3\sqrt{3}$.}
\end{ex}

% nguồn: Ok_TuanLuan_Toan11_Lan2
\begin{ex}%[1H8H7-5]%[ma:21133546]
	Cho hình lăng trụ  tứ giác đều có cạnh đáy bằng $a\sqrt{2}$, cạnh bên bằng $a\sqrt{5}$. Thể tích của khối lăng trụ đó bằng
	 \loigiai{
		Lăng trụ đã cho  là lăng trụ tứ giác đều nên đáy là hình vuông cạnh bằng $a \sqrt{2}$. Cạnh bên vuông góc với mặt đáy.\\
		$\Rightarrow$ Diện tích đáy của hình  lăng trụ là $B=(a \sqrt{2})^2=2a^2$.\\
		Vậy thể tích của khối lăng trụ đã cho bằng $V=B\cdot h=2a^2 \cdot  a\sqrt{5}=2a^3\sqrt{5}$.
	 }
\end{ex}

% nguồn: Ok_TuanLuan_Toan11_Lan2
\begin{ex}%[1H8V7-5]%[ma:f8780df3]
	Cho khối hộp  hình chữ nhật $ABCD.A'B'C'D'$ có đáy hình vuông, $AC=2\sqrt{3}a$,\break $\widehat{\left(\left(C'BD\right),\left(ABCD\right)\right)}=60^{\circ}$. Thể tích của khối hộp chữ nhật đã cho bằng
	\loigiai{ 
	\immini{ Gọi $O=AC\cap BD\Rightarrow OC=\dfrac{AC}{2}=a\sqrt{3}$,  $AB=\dfrac{AC}{\sqrt{2}}=a\sqrt{6}$.\\
		Ta có $\heva{&BD=\left(C'BD\right)\cap\left(ABCD\right)\\ 
		 &BD\perp\left(ACC'A'\right)\\
		 &OC'=\left(ACC'A'\right)\cap\left(ABCD\right)\\
		&OC=\left(ACC'A'\right)\cap\left(C'BD\right)}$  \\
		$\Rightarrow\widehat{\left(\left(C'BD\right),\left(ABCD\right)\right)}=\widehat{\left(OC',OC\right)}=\widehat{COC'}=60^{\circ}\left(\widehat{COC'} < 90^{\circ}\right)$.\\ 
		Xét tam giác $COC'$ vuông tại $C$.\\ 
		Ta có $\tan\widehat{COC'}=\dfrac{CC'}{OC}\Leftrightarrow CC'=OC\cdot \tan\widehat{COC'}=a\sqrt{3}\cdot \tan 60^{\circ}=3a$.  \\
		Ta có  $V_{ABCDA'B'C'D'}=S_{ABCD}\cdot CC'=\left(a\sqrt{6}\right)^2\cdot 3a=18a^3$.
		}{		\begin{tikzpicture}[scale=0.6,font=\footnotesize,line join=round,line cap=round,>=stealth]
		\def\bc{4} 
		\def\ba{2.5} 
		\def\gocB{65} 
		\coordinate[label=below left:$B$]  (B) at (0,0);
		 \coordinate[label=left:$A$] (A) at (\gocB:\ba);
		\coordinate[label=below:$C$]  (C) at (\bc,0);
		\coordinate[label=right:$D$]  (D) at ($(C)-(B)+(A)$);
		\coordinate[label=above  left:$A'$] (A') at ($(A)+(90:\bc)$);
		 \coordinate[label=left:$B'$] (B') at ($(B)-(A)+(A')$);
		\coordinate[label=right:$C'$] (C')  at ($(C)-(A)+(A')$);
		 \coordinate[label=right:$D'$] (D') at ($(D)-(A)+(A')$);
		\coordinate [label=left:$O$]  (O) at (intersection of A--C and B--D);
		\draw  (B')--(B)--(C)--(D)--(D')--(A')--(B')--(C')--(D')
			 (C)--(C');
		\draw[dashed]  (A)--(A')
			(A)--(B) 
			 (A)--(D)
			 (A)--(C)
			(B)--(D) 
			(O)--(C'); 
		 \draw pic["$60^\circ$",draw,angle eccentricity=2,angle radius=0.3cm]{angle=C--O--C'};
		\foreach \diem in {A,B,C,D,A',B',C',D',O}\fill  (\diem) circle (.5pt);
	\end{tikzpicture}
		 }
	 }
\end{ex}

% nguồn: Ok_TuanLuan_Toan11_Lan2
\begin{ex}%[1H8V7-5]%[ma:10185f2c]
	Cho hình hộp đứng $ ABCD.A'B'C'D'$ có đáy  là hình vuông cạnh $a$, góc giữa mặt phẳng $\left(D'AB\right)$ và mặt phẳng $\left(ABCD\right)$ là $30^\circ$. Thể tích khối hộp $ ABCD.A'B'C'D'$ bằng
	 \loigiai{
	\immini{ Ta  có $AB\perp\left(ADD'A'\right)\Rightarrow AB\perp D'A$.\\
		Lại có $ AB\perp AD$.\\ 
		Suy  ra $\widehat{\left(\left(D'AB\right);\left(ABCD\right)\right)}=\widehat{\left(D'A,AD\right)}=\widehat{D'AD}=30^\circ$.
		 \\
		Xét $\triangle  D'DA$ vuông tại $D$ có $AD=a$.\\
		Suy ra  $\widehat{D'AD}=30^\circ$ $\Rightarrow DD'=\tan{30^\circ}\cdot AD=\dfrac{a\sqrt{3}}{3}$.\\
		Vậy $V_{ABCD.A'B'C'D'}=DD'\cdot  S_{ABCD}=\dfrac{a\sqrt{3}}{3}\cdot a^2=\dfrac{a^3\sqrt{3}}{3}$.
		}{		\begin{tikzpicture}[scale=0.6,font=\footnotesize,line join=round,line cap=round,>=stealth]
		\def\bc{4} 
		\def\ba{2.5} 
		 \def\gocB{65}
		\coordinate[label=below  left:$B$] (B) at (0,0);
		\coordinate[label=left:$A$] (A)  at (\gocB:\ba);
		\coordinate[label=below:$C$] (C) at  (\bc,0);
		\coordinate[label=right:$D$]  (D) at ($(C)-(B)+(A)$);
		\coordinate[label=above left:$A'$] (A')  at ($(A)+(90:\bc)$);
		\coordinate[label=left:$B'$] (B') at  ($(B)-(A)+(A')$);
		\coordinate[label=below  right:$C'$] (C') at ($(C)-(A)+(A')$);
		 \coordinate[label=right:$D'$] (D') at ($(D)-(A)+(A')$);
		\draw (B')--(B)--(C)--(D)--(D')--(A')--(B')--(C')--(D') 
			(C)--(C'); 
		\draw[dashed] (A)--(D')--(B) 
			 (A)--(A')
			 (A)--(B)
			 (A)--(D);
		\foreach \diem  in {A,B,C,D,A',B',C',D'} \fill (\diem) circle (.5pt);
	\end{tikzpicture}
		 }
	 }
\end{ex}

% nguồn: Ok_TuanLuan_Toan11_Lan2
\begin{ex}%[1H8C7-5]%[ma:96b42dab]
	Cho hình lập phương $ABCD.A'B'C'D'$ có khoảng cách giữa hai đường thẳng $AB'$ và $BD$ bằng $\dfrac{2\sqrt{3}a}{3}$. Thể tích của khối  lập phương $ABCD.A'B'C'D'$ bằng
	\loigiai{ 
	\immini{ Gọi  $O$ là giao điểm của $BD$ và $AC$.\\
		Ta có $\heva{& BD\perp AC\\&  BD\perp C{C}'\\& AC\cap C{C}'=C}
		\Rightarrow  BD\perp\left(AC{C}'A'\right)$.\\
		Trong $\left(AC{C}'A'\right)$: Từ $C$ hạ $CH\perp C'O$  tại $H$ \\
		Khi đó ta có  $\heva{& CH\perp BD\\& CH\perp{C}'O\\& C'O\cap BD=O}\Rightarrow CH\perp\left(BDC'\right)$ \\
		Ta lại  có
		$\heva{&AB'\parallel DC'\subset\left(BDC'\right)\\&AB'\not\subset\left(BDC'\right)} 
		 \Rightarrow AB'\parallel\left(BDC'\right)$
		}{		\begin{tikzpicture}[scale=0.7,font=\footnotesize,line join=round,line cap=round,>=stealth]
		\def\bc{6} 
		 \def\ba{2.5}
		\def\gocB{150} 
		 \coordinate[label=below left:$B$] (B) at (0,0);
		 \coordinate[label=left:$A$] (A) at (\gocB:\ba);
		 \coordinate[label=below:$C$] (C) at (\bc,0);
		\coordinate[label=below:$D$] (D) at ($(C)-(B)+(A)$); 
		\coordinate[label=above left:$A'$]  (A') at ($(A)+(90:\bc)$);
		\coordinate[label=above:$B'$] (B')  at ($(B)-(A)+(A')$);
		\coordinate[label=right:$C'$] (C') at ($(C)-(A)+(A')$); 
		\coordinate[label=above:$D'$] (D')  at ($(D)-(A)+(A')$);
		\coordinate [label=below:$O$] (O)  at (intersection of A--C and B--D);
		\coordinate[label=left:$H$] (H) at ($(O)!(C)!(C')$); 
		\draw pic[draw,angle radius=0.25cm]{right  angle=O--H--C};
		\draw (B')--(B) 
			(A')--(B') 
			 (B)--(C)
			(D')--(A') 
			 (B')--(C')--(D')
			 (C)--(C')
			(B)--(C') 
			 (A)--(B)
			(A)--(A') 
			(A)--(B'); 
		 \draw[dashed] (A)--(D)--(D')
			 (D)--(C)
			 (A)--(C)
			(B)--(D)--(C')--(O) 
			 (C)--(H);
		\foreach \diem in {A,B,C,D,A',B',C',D',O,H}\fill (\diem) circle (1pt); 
	\end{tikzpicture}
		} 
		Khi đó, ta có 
		\begin{eqnarray*}
			\mathrm{d}\left(AB';BD\right)&=&\mathrm{d}\left(AB';\left(BDC'\right)\right)\\ 
			&=&\mathrm{d}\left(A;\left(BDC'\right)\right)\\ 
			&=&\mathrm{d}\left(C,\left(BDC'\right)\right)\\ 
			&=&CH\\ 
			&=&\dfrac{2a\sqrt{3}}{3} 
		\end{eqnarray*}
		Đặt cạnh  hình lập phương là $x$ $\Rightarrow\heva{& CC'=x\\ & CO=\dfrac{x\sqrt{2}}{2}}$.\\
		Khi đó $\dfrac{1}{CH^2}=\dfrac{1}{CC'^2}+\dfrac{1}{CO^2} 
		 \Leftrightarrow\dfrac{3}{4a^2}=\dfrac{1}{x^2}+\dfrac{2}{x^2}
		=\dfrac{3}{x^2}\Leftrightarrow{x^2}=4a^2\Leftrightarrow x=2a$.\\ 
		Do đó thể tích  của khối lập phương là $(2a)^3=8a^3$.
	} 
\end{ex}
